\ifdefined\gkver\else\def\gkver{student}\fi
\documentclass[\gkver]{gaokaozhenti}
\begin{document}

\chapter{2025年云南}

\begin{enumerate}
\item
%% number: 1
%% paperName: 2025年云南高考第 1 题 4 分
%% typeId: 1
%% type: 单选题
%% chapter: 原子和原子核
%% point: 天然放射性
%% method: 无
%% score: 4
%% degree: 650
%% duplicateId: 0
%% body:
2025 年 3 月，我国科学家研制的碳 14 核电池原型机“烛龙一号”发布，标志着我国在核能技术领域与微型核电池领域取得突破。碳 14 的衰变方程为 \ce{^{14}_{6} C -> ^{14}_{7} N} + X，则\xzanswer{A}

\fourchoices[answer=A]
{X 为电子，是在核内中子转化为质子的过程中产生的}
{X 为电子，是在核内质子转化为中子的过程中产生的}
{X 为质子，是由核内中子转化而来的}
{X 为中子，是由核内质子转化而来的}

%% answer: A
%% memo:
\memoanswer{
根据质量数和电荷数守恒有 \ce{^{14}_{6} C -> ^{14}_{7} N + ^{0}_{-1} e}

可知 X 为电子，电子是在核内中子转化为质子的过程中产生的。

故选 A。
}

\item
%% number: 2
%% paperName: 2025年云南高考第 2 题 4 分
%% typeId: 1
%% type: 单选题
%% chapter: 功和能
%% point: 动能定理
%% method: 无
%% score: 4
%% degree: 850
%% duplicateId: 0
%% body:
如图所示，中老铁路国际旅客列车从云南某车站由静止出发，沿水平直轨道逐渐加速到 $144 \Ukmh$，在此过程中列车对座椅上的一高中生所做的功最接近\xzanswer{B}

\onepicture[w=5.10cm, h=2.87cm]{figs/02.jpg}

\fourchoices[answer=B]
{$4 \times 10^{5} \UJ$}
{$4 \times 10^{4} \UJ$}
{$4 \times 10^{3} \UJ$}
{$4 \times 10^{2} \UJ$}

%% answer: B
%% memo:
\memoanswer{
高中生的质量约为 $50 \Ukg$，根据动能定理有 $W = \frac{1}{2}mv^{2} = 4 \times 10^{4} \UJ$。

故选 B。
}

\item
%% number: 3
%% paperName: 2025年云南高考第 3 题 4 分
%% typeId: 1
%% type: 单选题
%% chapter: 曲线运动
%% point: 平抛运动
%% method: 无
%% score: 4
%% degree: 550
%% duplicateId: 0
%% body:
如图所示，某同学将两颗鸟食从 O 点水平抛出，两只小鸟分别在空中的 M 点和 N 点同时接到鸟食。鸟食的运动视为平抛运动，两运动轨迹在同一竖直平面内，则\xzanswer{D}

\onepicture[w=5.00cm, h=2.77cm]{figs/03.png}

\fourchoices[answer=D]
{两颗鸟食同时抛出}
{在 N 点接到的鸟食后抛出}
{两颗鸟食平抛的初速度相同}
{在 M 点接到的鸟食平抛的初速度较大}

%% answer: D
%% memo:
\memoanswer{
AB．鸟食的运动视为平抛运动，则在竖直方向有 $h = \frac{1}{2}gt^{2}$。

\onepicture[w=4.90cm, h=2.73cm]{figs/03_an.png}

由于 $h_{M} < h_{N}$，则 $t_{M} < t_{N}$，要同时接到鸟食，则在 N 点接到的鸟食先抛出，故 AB 错误；

CD．在水平方向有 $x = v_{0}t$，如图所示，过 M 点作一水平面，可看出在相同高度处 M 点的水平位移大，则 M 点接到的鸟食平抛的初速度较大，故 C 错误，D 正确。

故选 D。
}

\item
%% number: 4
%% paperName: 2025年云南高考第 4 题 4 分
%% typeId: 1
%% type: 单选题
%% chapter: 电场
%% point: 等势面
%% method: 无
%% score: 4
%% degree: 550
%% duplicateId: 0
%% body:
某介电电泳实验使用非匀强电场，该电场的等势线分布如图所示。a、b、c、d 四点分别位于电势为 $-2 \UV$、$-1 \UV$、$1 \UV$、$2 \UV$ 的等势线上，则\xzanswer{C}

\onepicture[w=4.20cm, h=4.03cm]{\includesvg{figs/04}}

\fourchoices[answer=C]
{a、b、c、d 中 a 点电场强度最小}
{a、b、c、d 中 d 点电场强度最大}
{一个电子从 b 点移动到 c 点电场力做功为 $2 \UeV$}
{一个电子从 a 点移动到 d 点电势能增加了 $4 \UeV$}

%% answer: C
%% memo:
\memoanswer{
AB．根据等势面越密集电场强度越大，可知 a、b、c、d 中 a 点电场强度最大，故 AB 错误；

C．一个电子从 b 点移动到 c 点电场力做功为 $W_{bc} = -eU_{bc} = 2 \UeV$

故 C 正确；

D．一个电子从 a 点移动到 d 点电场力做功为 $W_{ad} = -eU_{bc} = 4 \UeV$

由于电场力做正功电势能减小，则一个电子从 a 点移动到 d 点电势能减小了 $4 \UeV$，故 D 错误。

故选 C。
}

\item
%% number: 5
%% paperName: 2025年云南高考第 5 题 4 分
%% typeId: 1
%% type: 单选题
%% chapter: 曲线运动
%% point: 开普勒定律
%% method: 无
%% score: 4
%% degree: 550
%% duplicateId: 0
%% body:
国际编号为 192391 的小行星绕太阳公转的周期约为 5.8 年，该小行星与太阳系内八大行星几乎在同一平面内做圆周运动。规定地球绕太阳公转的轨道半径为 1 AU，八大行星绕太阳的公转轨道半径如下表所示。忽略其它行星对该小行星的引力作用，则该小行星的公转轨道应介于\xzanswer{C}

\begin{center}
\small
\begin{tabular}{|c|c|c|c|c|c|c|c|c|}
\hline
行星 & 水星 & 金星 & 地球 & 火星 & 木星 & 土星 & 天王星 & 海王星 \\
\hline
轨道半径 $R/\UAU$ & 0.39 & 0.72 & 1.0 & 1.5 & 5.2 & 9.5 & 19 & 30 \\
\hline
\end{tabular}
\end{center}

\fourchoices[answer=C]
{金星与地球的公转轨道之间}
{地球与火星的公转轨道之间}
{火星与木星的公转轨道之间}
{天王星与海王星的公转轨道之间}

%% answer: C
%% memo:
\memoanswer{
根据开普勒第三定律可知 $\frac{r_{\text{行}}^{3}}{T_{\text{行}}^{2}} = \frac{r_{\text{地}}^{3}}{T_{\text{地}}^{2}}$。

其中 $r_{\text{地}} = 1 \UAU$，$T_{\text{地}} = 1$ 年，$T_{\text{行}} = 5.8$ 年

代入解得 $r_{\text{行}} \approx 3.23 \UAU$

故可知该小行星的公转轨道应介于火星与木星的公转轨道之间。

故选 C。
}

\item
%% number: 6
%% paperName: 2025年云南高考第 6 题 4 分
%% typeId: 1
%% type: 单选题
%% chapter: 功和能
%% point: 动能定理
%% method: 无
%% score: 4
%% degree: 450
%% duplicateId: 0
%% body:
如图所示，质量为 $m$ 的滑块（视为质点）与水平面上 MN 段的动摩擦因数为 $\mu_{1}$，与其余部分的动摩擦因数为 $\mu_{2}$，且 $\mu_{1} > \mu_{2}$。第一次，滑块从 I 位置以速度 $v_{0}$ 向右滑动，通过 MN 段后停在水平面上的某一位置，整个运动过程中，滑块的位移大小为 $x_{1}$，所用时间为 $t_{1}$；第二次，滑块从 Ⅱ 位置以相同速度 $v_{0}$ 向右滑动，通过 MN 段后停在水平面上的另一位置，整个运动过程中，滑块的位移大小为 $x_{2}$，所用时间为 $t_{2}$。忽略空气阻力，则\xzanswer{A}

\onepicture[w=6.50cm, h=1.74cm]{\includesvg{figs/06}}

\fourchoices[answer=A]
{$t_{1} < t_{2}$}
{$t_{1} > t_{2}$}
{$x_{1} > x_{2}$}
{$x_{1} < x_{2}$}

%% answer: A
%% memo:
\memoanswer{
CD．对两种运动的整个过程根据能量守恒有

$\frac{1}{2}mv_{0}^{2} = \mu_{1}mgx_{MN} + \mu_{2}mg(x_{1} - x_{MN})$，$\frac{1}{2}mv_{0}^{2} = \mu_{1}mgx_{MN} + \mu_{2}mg(x_{2} - x_{MN})$

可得 $x_{1} = x_{2}$，故 CD 错误；

AB．根据牛顿第二定律 $\mu mg = ma$

可得 $a = \mu g$

由于 $\mu_{1} > \mu_{2}$，故滑块在 MN 上时的加速度大，根据前面分析可知两次运动的总位移相等，即两次运动过程中 $v-t$ 图像与横轴围成的面积相等，由于第二次时滑块距离 M 点的距离较近，根据公式 $v_{0}^{2} - v^{2} = 2\mu_{2}gx$ 可知第二次到达 M 点时速度较大，作出整个过程中两种运动状态的 $v-t$ 图像可得 $t_{1} < t_{2}$，故 A 正确，B 错误；

\onepicture[w=4.60cm, h=3.68cm]{figs/06_an.png}

故选 A。
}

\item
%% number: 7
%% paperName: 2025年云南高考第 7 题 4 分
%% typeId: 1
%% type: 单选题
%% chapter: 机械波
%% point: 波长频率波速
%% method: 无
%% score: 4
%% degree: 450
%% duplicateId: 0
%% body:
如图所示，均匀介质中矩形区域内有一位置未知的波源。$t = 0$ 时刻，波源开始振动产生简谐横波，并以相同波速分别向左、右两侧传播，P、Q 分别为矩形区域左右两边界上振动质点的平衡位置。$t = 1.5 \Us$ 和 $t = 2.5 \Us$ 时矩形区域外波形分别如图中实线和虚线所示，则\xzanswer{D}

\onepicture[w=6.50cm, h=2.36cm]{figs/07.png}

\fourchoices[answer=D]
{波速为 $2.5 \Ums$}
{波源的平衡位置距离 P 点 $1.5 \Um$}
{$t = 1.0 \Us$ 时，波源处于平衡位置且向下运动}
{$t = 5.5 \Us$ 时，平衡位置在 P、Q 处的两质点位移相同}

%% answer: D
%% memo:
\memoanswer{
A．根据波形可知 $\lambda = 4 \Um$，$T/2 = 2.5 \Us - 1.5 \Us$，可得 $T = 2 \Us$。故波速为 $v = \frac{\lambda}{T} = 2 \Ums$。

故 A 错误；

B．设波源的平衡位置距离 P 点距离为 $x_{0}$，根据左侧 $t = 1.5 \Us$ 时的波形可知 $\frac{2 + x_{0}}{v} = 1.5 \Us$

解得 $x_{0} = 1 \Um$

故 B 错误；

C．根据左侧实线波形结合同侧法可知波源刚开始的振动方向向下，由于 $t = 1.0 \Us = T/2$，故可知此时波源处于平衡位置且向上运动，故 C 错误；

D．由于 $x_{0} = 1 \Um$，可知波源的平衡位置距离 Q 点距离为 $x_{1} = 3 \Um$。

故波传到 P、Q 两点的时间分别为 $t_{0} = \frac{x_{0}}{v} = 0.5 \Us$，$t_{1} = \frac{x_{1}}{v} = 1.5 \Us$

故 $t = 5.5 \Us$ 时，平衡位置在 P、Q 处的两质点已经振动的时间分别为 $t_{0}' = 5.5 \Us - 0.5 \Us = \frac{5}{2}T$，$t_{1}' = 5.5 \Us - 1.5 \Us = 2T$

由于波源刚开始向下振动，故 $t = 5.5 \Us$ 时，P 处质点处于平衡位置向上振动，Q 处质点处于平衡位置向下振动，故此时平衡位置在 P、Q 处的两质点位移相同。

故 D 正确。

故选 D。
}

\item
%% number: 8
%% paperName: 2025年云南高考第 8 题 6 分
%% typeId: 2
%% type: 多选题
%% chapter: 交流电
%% point: 变压器
%% method: 无
%% score: 6
%% degree: 550
%% duplicateId: 0
%% body:
电动汽车充电桩的供电变压器（视为理想变压器）示意图如图所示。变压器原线圈的匝数为 $n_{1}$，输入电压 $U_{1} = 1.1 \UkV$；两副线圈的匝数分别为 $n_{2}$ 和 $n_{3}$，输出电压 $U_{2} = U_{3} = 220 \UV$。当 I、Ⅱ 区充电桩同时工作时，两副线圈的输出功率分别为 $7.0 \UkW$ 和 $3.5 \UkW$，下列说法正确的是\xzanswer{AC}

\onepicture[w=6.00cm, h=3.73cm]{figs/08.png}

\fourchoices[answer=AC]
{$n_{1}:n_{2} = 5:1$}
{$n_{1}:n_{3} = 1:5$}
{变压器的输入功率为 $10.5 \UkW$}
{两副线圈输出电压最大值均为 $220 \UV$}

%% answer: AC
%% memo:
\memoanswer{
AB．根据理想变压器的电压比等于匝数比可得 $n_{1}:n_{2} = U_{1}:U_{2} = 5:1$，$n_{1}:n_{3} = U_{1}:U_{3} = 5:1$。

故 A 正确，B 错误；

C．根据能量守恒可知变压器的输入功率等于总的输出功率，故

$P_{\text{输入}} = P_{\text{输出}} = (7.0 + 3.5) \UkW = 10.5 \UkW$

故 C 正确；

D．输出电压为交流电的有效值，根据正弦交流电的最大值与有效值的关系可知，两副线圈输出电压最大值均为 $U_{m} = 220\sqrt{2} \UV$

故 D 错误。

故选 AC。
}

\item
%% number: 9
%% paperName: 2025年云南高考第 9 题 4 分
%% typeId: 2
%% type: 多选题
%% chapter: 气体性质
%% point: 查理定律
%% method: 无
%% score: 4
%% degree: 550
%% duplicateId: 0
%% body:
图~\ref{2025云南09a}为 1593 年伽利略发明的人类历史上第一支温度计，其原理如图~\ref{2025云南09b}所示。硬质玻璃泡 a 内封有一定质量的气体（视为理想气体），与 a 相连的 b 管插在水槽中固定，b 管中液面高度会随环境温度变化而变化。设 b 管的体积与 a 泡的体积相比可忽略不计，在标准大气压 $p_{0}$ 下，b 管上的刻度可以直接读出环境温度。则在 $p_{0}$ 下\xzanswer{BD}

\twopicture[labelA=2025云南09a, labelB=2025云南09b, h=3.30cm]{figs/09a.png}{figs/09b.png}

\fourchoices[answer=BD]
{环境温度升高时，b 管中液面升高}
{环境温度降低时，b 管中液面升高}
{水槽中的水少量蒸发后，温度测量值偏小}
{水槽中的水少量蒸发后，温度测量值偏大}

%% answer: BD
%% memo:
\memoanswer{
AB．根据题意，a 中气体做等容变化，根据查理定律可知，当环境温度升高，a 中气体压强增大，又 $p_{a} + \rho gh = p_{0}$，可知 b 管中液面降低，同理可知环境温度降低时，b 管中液面升高，故 B 正确，A 错误；

CD．由 AB 选项分析可知，b 管中刻度从上到下温度逐渐升高，同一温度，a 中压强不变，b 管中液面液槽内液面高度差不变，水槽中的水少量蒸发后，槽中液面降低，则 b 管内液面降低，则温度测量值偏大，故 D 正确，C 错误。

故选 BD。
}

\item
%% number: 10
%% paperName: 2025年云南高考第 10 题 6 分
%% typeId: 1
%% type: 单选题
%% chapter: 功和能
%% point: 动能定理
%% method: 无
%% score: 6
%% degree: 350
%% duplicateId: 0
%% body:
如图所示，倾角为 $\theta$ 的固定斜面，其顶端固定一劲度系数为 $k$ 的轻质弹簧，弹簧处于原长时下端位于 O 点。质量为 $m$ 的滑块 Q（视为质点）与斜面间的动摩擦因数 $\mu = \tan \theta$。过程 Ⅰ：Q 以速度 $v_{0}$ 从斜面底端 P 点沿斜面向上运动恰好能滑至 O 点；过程 Ⅱ：将 Q 连接在弹簧的下端并拉至 P 点由静止释放，Q 通过 M 点（图中未画出）时速度最大，过 O 点后能继续上滑。弹簧始终在弹性限度内，假设最大静摩擦力等于滑动摩擦力，忽略空气阻力，重力加速度为 $g$。则\xzanswer{BCD}

\onepicture[w=5.30cm, h=3.63cm]{figs/10.png}

\fourchoices[answer=BCD]
{P、M 两点之间的距离为 $\frac{kv_{0}^{2} - 4mg^{2}\sin^{2}\theta}{4kg\sin\theta}$}
{过程 Ⅱ 中，Q 在从 P 点单向运动到 O 点的过程中损失的机械能为 $\frac{1}{4}mv_{0}^{2}$}
{过程 Ⅱ 中，Q 从 P 点沿斜面向上运动的最大位移为 $\frac{kv_{0}^{2} - 8mg^{2}\sin^{2}\theta}{2kg\sin\theta}$}
{连接在弹簧下端的 Q 无论从斜面上何处释放，最终一定静止在 OM（含 O、M 点）之间}

%% answer: BCD
%% memo:
\memoanswer{
A．设 PO 的距离为 $L$，过程 Ⅰ，根据动能定理有

$-mg \sin \theta \cdot L - \mu mg \cos \theta \cdot L = 0 - \frac{1}{2}mv_{0}^{2}$

设 MO 的距离为 $L_{1}$，过程 Ⅱ 中，当 Q 速度最大时，根据平衡条件

$kL_{1} = mg \sin \theta + \mu mg \cos \theta$

P、M 两点之间的距离 $L_{2} = L - L_{1}$

联立可得 $L_{2} = \frac{kv_{0}^{2} - 8mg^{2}\sin^{2}\theta}{4kg\sin\theta}$

故 A 错误；

B．根据功能关系，可知过程 Ⅱ 中，Q 在从 P 点单向运动到 O 点的过程中损失的机械能

$\Delta E = \mu mg \cos \theta \cdot L$

结合 $-mg \sin \theta \cdot L - \mu mg \cos \theta \cdot L = 0 - \frac{1}{2}mv_{0}^{2}$

可得 $\Delta E = \frac{1}{4}mv_{0}^{2}$

故 B 正确；

C．设过程 Ⅱ 中，Q 从 P 点沿斜面向上运动的最大位移为 $x$，根据能量守恒定律

$\frac{1}{2}kL^{2} = mg \sin \theta \cdot x + \mu mg \cos \theta \cdot x + \frac{1}{2}k(x - L)^{2}$

结合 $-mg \sin \theta \cdot L - \mu mg \cos \theta \cdot L = 0 - \frac{1}{2}mv_{0}^{2}$

解得 $x = \frac{kv_{0}^{2} - 8mg^{2}\sin^{2}\theta}{2kg\sin\theta}$

故 C 正确；

D．无论 Q 从何处释放，Q 在斜面上运动过程中，弹簧与 Q 初始时的势能变为摩擦热，当在 M 点时，满足 $kL_{1} = mg \sin \theta + \mu mg \cos \theta$

当在 O 点时，满足 $mg \sin \theta = \mu mg \cos \theta$

所以在 OM（含 O、M 点）之间速度为零时，Q 将静止，故 D 正确。

故选 BCD。
}

\item
%% number: 11
%% paperName: 2025年云南高考第 11 题 6 分
%% typeId: 4
%% type: 实验题
%% chapter: 力物体的平衡
%% point: 摩擦力
%% method: 无
%% score: 6
%% degree: 650
%% duplicateId: 0
%% body:
某实验小组做了测量木质滑块与橡胶皮之间动摩擦因数 $\mu$ 的实验，所用器材如下：钉有橡胶皮的长木板、质量为 $250 \Ug$ 的木质滑块（含挂钩）、细线、定滑轮、弹簧测力计、慢速电机以及砝码若干。实验装置如图~\ref{2025云南11a}所示。

实验步骤如下：
\begin{steps}
\item
将长木板放置在水平台面上，滑块平放在橡胶面上；
\item
调节定滑轮高度，使细线与长木板平行（此时定滑轮高度与挂钩高度一致）；
\item
用电机缓慢拉动长木板，当长木板相对滑块匀速运动时，记录弹簧测力计的示数 $F$；
\item
在滑块上分别放置 $50 \Ug$、$100 \Ug$ 和 $150 \Ug$ 的砝码，重复步骤③；
\item
处理实验数据（重力加速度 $g$ 取 $9.80 \Umsq$）。
\end{steps}

实验数据如下表所示：
\begin{center}
\begin{tabular}{|c|c|c|}
\hline
滑块和砝码的总质量 $M/\Ug$ & 弹簧测力计示数 $F/\UN$ & 动摩擦因数 $\mu$ \\
\hline
250 & 1.12 & 0.457 \\
\hline
300 & 1.35 & a \\
\hline
350 & 1.57 & 0.458 \\
\hline
400 & 1.79 & 0.457 \\
\hline
\end{tabular}
\end{center}

完成下列填空：
\begin{enumerate}
\item
表格中 a 处的数据为 \tkanswer{0.459}（保留 3 位有效数字）；
\item
其它条件不变时，在实验误差允许的范围内，滑动摩擦力的大小与接触面上压力的大小 \tkanswer{成正比}，$\mu$ 与接触面上压力的大小 \tkanswer{无关}（以上两空填“成正比”“成反比”或“无关”）；
\item
若在实验过程中未进行步骤②，实验装置如图~\ref{2025云南11b}所示，挂钩高于定滑轮，则 $\mu$ 的测量结果将 \tkanswer{偏大}（填“偏大”“偏小”或“不变”）。
\end{enumerate}

\twopicture[labelA=2025云南11a, labelB=2025云南11b, w=6.50cm, align=right]{figs/11a.png}{figs/11b.png}

%% answer:
\jdanswer{
\begin{enumerate}
\item
0.459

\item
成正比，无关

\item
偏大
\end{enumerate}
}
%% memo:
\memoanswer{
（1）表格中 a 处的数据 $\mu = \frac{1.35}{0.3 \times 9.8} \approx 0.459$

（2）根据表中数据分析可知其它条件不变时，在实验误差允许的范围内，滑动摩擦力的大小与接触面上压力的大小成正比；

根据表中数据分析可知其它条件不变时，在实验误差允许的范围内，$\mu$ 与接触面上压力的大小无关。

（3）实验装置中挂钩高于定滑轮，则绳子拉力有竖直向下的分力，实际的正压力大于测量值的正压力，即 $F_{\text{压测}} < F_{\text{压实}}$

根据 $\mu = \frac{F}{F_{\text{压测}}}$

可得 $\mu_{\text{测}} > \mu_{\text{实}}$
}

\item
%% number: 12
%% paperName: 2025年云南高考第 12 题 10 分
%% typeId: 4
%% type: 实验题
%% chapter: 电路
%% point: 传感器
%% method: 无
%% score: 10
%% degree: 550
%% duplicateId: 0
%% body:
基于铂电阻阻值随温度变化的特性，某兴趣小组用铂电阻做了测量温度的实验。可选用的器材如下：Pt1000 型号铂电阻、电源 $E$（电动势 $5 \UV$，内阻不计）、电流表 A$_{1}$（量程 $100 \UuA$，内阻 $4.5 \UkO$）、电流表 A$_{2}$（量程 $500 \UuA$，内阻约 $1 \UkO$）、定值电阻 $R_{1}$（阻值 $15 \UkO$）、定值电阻 $R_{2}$（阻值 $1.5 \UkO$）、开关 S 和导线若干。

查阅技术手册可知，Pt1000 型号铂电阻测温时的工作电流在 $0.1 \sim 0.3 \UmA$ 之间，在 $0 \sim 100\UdC$ 范围内，铂电阻的阻值 $R_{t}$ 随温度 $t$ 的变化视为线性关系，如图~\ref{2025云南12a}所示。

完成下列填空：
\begin{enumerate}
\item
由图~\ref{2025云南12a}可知，在 $0 \sim 100\UdC$ 范围内，温度每升高 $1\UdC$，该铂电阻的阻值增加 \tkanswer{3.85} $\UO$；
\item
兴趣小组设计了如图~\ref{2025云南12b}所示的甲、乙两种测量铂电阻阻值的电路图，能准确测出铂电阻阻值的是 \tkanswer{乙}（填“甲”或“乙”），保护电阻 $R$ 应选 \tkanswer{$R_{1}$}（填“$R_{1}$”或“$R_{2}$”）；
\item
用（2）问中能准确测出铂电阻阻值的电路测温时，某次测量读得 A$_{2}$ 示数为 $295 \UuA$，A$_{1}$ 示数如图~\ref{2025云南12c}所示，该示数为 \tkanswer{62.0} $\UuA$，则所测温度为 \tkanswer{51} $\UdC$（计算结果保留 2 位有效数字）。
\end{enumerate}

\threepicture[labelA=2025云南12a, labelB=2025云南12b, labelC=2025云南12c, wA=4.30cm, wB=9.00cm, wC=8.00cm, align=right]{figs/12a.png}{figs/12b.png}{figs/12c.png}

%% answer:
\jdanswer{
\begin{enumerate}
\item
3.85

\item
乙，$R_{1}$

\item
62.0，51
\end{enumerate}
}
%% memo:
\memoanswer{
（1）温度每升高 $1\UdC$，该铂电阻的阻值增加

$\Delta R = \frac{1.385 - 1.000}{100} \times 10^{3} \UO = 3.85 \UO$

（2）由于 A$_{1}$ 内阻确定，所以用 A$_{1}$ 测量电阻的电压，用 A$_{2}$ 与 A$_{1}$ 之差来测量经过电阻的电流，故能准确测出铂电阻阻值的是乙；

电路中的最大电流为 $0.3 \UmA$，可得电路中的最小阻值

$R_{\min} = \frac{5}{0.0003} \UO \approx 17 \UkO$

可知保护电阻 $R$ 应选 $R_{1}$。

（3）由图可知 A$_{1}$ 的分度值为 $1 \UuA$，则其读数为 $62.0 \UuA$；

根据欧姆定律可得

$R = \frac{I_{1}R_{A1}}{I_{2} - I_{1}}$

根据题图可得 $R = 1000 + 3.85t$

代入数据可得 $t \approx 51\UdC$
}

\item
%% number: 13
%% paperName: 2025年云南高考第 13 题 10 分
%% typeId: 6
%% type: 计算题
%% chapter: 光学
%% point: 全反射
%% method: 无
%% score: 10
%% degree: 450
%% duplicateId: 0
%% body:
用光学显微镜观察样品时，显微镜部分结构示意图如图~\ref{2025云南13a}所示。盖玻片底部中心位置 O 点的样品等效为点光源，为避免 O 点发出的光在盖玻片上方界面发生全反射，可将盖玻片与物镜的间隙用一滴油填充，如图~\ref{2025云南13b}所示。已知盖玻片材料和油的折射率均为 1.5，盖玻片厚度 $d = 2.0 \Umm$，盖玻片与物镜的间距 $h = 0.20 \Umm$，不考虑光在盖玻片中的多次反射，取真空中光速 $c = 3.0 \times 10^{8} \Ums$，$\pi = 3.14$。

\begin{enumerate}
\item
求未滴油时，O 点发出的光在盖玻片的上表面的透光面积（结果保留 2 位有效数字）；
\item
滴油前后，光从 O 点传播到物镜的最短时间分别为 $t_{1}$、$t_{2}$，求 $t_{2} - t_{1}$（结果保留 2 位有效数字）。
\end{enumerate}

\twopicture[labelA=2025云南13a, labelB=2025云南13b, capA=填充前, capB=填充后, w=6.30cm, align=right]{figs/13a.png}{figs/13b.png}

%% answer:
\jdanswer{
\begin{enumerate}
\item
$1.0 \times 10^{-5} \Umq$

\item
$3.3 \times 10^{-13} \Us$
\end{enumerate}
}
%% memo:
\memoanswer{
（1）由折射定律可知，全反射的临界角满足 $\sin C = \frac{1}{n} = \frac{2}{3}$

设未滴油时，O 点发出的光在盖玻片的上表面的透光圆的半径为 $r$，由几何关系

$\sin C = \frac{r}{\sqrt{r^{2} + d^{2}}}$

代入数据解得 $r = \frac{4\sqrt{5}}{5} \Umm$

根据 $S = \pi r^{2}$

所以未滴油时，O 点发出的光在盖玻片的上表面的透光面积为 $S \approx 1.0 \times 10^{-5} \Umq$

（2）当光从 O 点垂直于盖玻片的上表面入射时，传播的时间最短，则未滴油时，光从 O 点传播到物镜的最短时间为

$t_{1} = \frac{d}{v} + \frac{h}{c} = \frac{d}{\frac{c}{n}} + \frac{h}{c} = \frac{nd + h}{c}$

滴油时，光从 O 点传播到物镜的最短时间为

$t_{2} = \frac{d}{v} + \frac{h}{v} = \frac{d}{\frac{c}{n}} + \frac{h}{\frac{c}{n}} = \frac{n(d + h)}{c}$

故 $t_{2} - t_{1} = \frac{(n - 1)h}{c} = \frac{0.5 \times 0.2 \times 10^{-3}}{3.0 \times 10^{8}} \Us \approx 3.3 \times 10^{-13} \Us$
}

\item
%% number: 14
%% paperName: 2025年云南高考第 14 题 13 分
%% typeId: 6
%% type: 计算题
%% chapter: 磁场
%% point: 带电粒子在磁场中的运动
%% method: 无
%% score: 13
%% degree: 450
%% duplicateId: 0
%% body:
磁屏蔽技术可以降低外界磁场对屏蔽区域的干扰。如图所示，$x \geq 0$ 区域存在垂直 $Oxy$ 平面向里的匀强磁场，其磁感应强度大小为 $B_{1}$（未知）。第一象限内存在边长为 $2L$ 的正方形磁屏蔽区 ONPQ，经磁屏蔽后，该区域内的匀强磁场方向仍垂直 $Oxy$ 平面向里，其磁感应强度大小为 $B_{2}$（未知），但满足 $0 < B_{2} < B_{1}$。某质量为 $m$、电荷量为 $q$（$q > 0$）的带电粒子通过速度选择器后，在 $Oxy$ 平面内垂直 $y$ 轴射入 $x \geq 0$ 区域，经磁场偏转后刚好从 ON 中点垂直 ON 射入磁屏蔽区域。速度选择器两极板间电压 $U$、间距 $d$、内部磁感应强度大小 $B_{0}$ 已知，不考虑该粒子的重力。

\begin{enumerate}
\item
求该粒子通过速度选择器的速率；
\item
求 $B_{1}$ 以及 $y$ 轴上可能检测到该粒子的范围；
\item
定义磁屏蔽效率 $\eta = \frac{B_{1} - B_{2}}{B_{1}} \times 100\%$，若在 Q 处检测到该粒子，则 $\eta$ 是多少？
\end{enumerate}

\onepicture[h=6.00cm, align=right]{figs/14.png}

%% answer:
\jdanswer{
\begin{enumerate}
\item
$\frac{U}{B_{0}d}$

\item
$\frac{mU}{qdB_{0}L}$，$L < y < 3L$

\item
$60\%$
\end{enumerate}
}
%% memo:
\memoanswer{
（1）由于该粒子在速度选择器中受力平衡，故 $qE = qv_{0}B_{0}$

其中 $E = \frac{U}{d}$

则该粒子通过速度选择器的速率为 $v_{0} = \frac{U}{B_{0}d}$

（2）粒子在 $x \geq 0$ 区域内做匀速圆周运动，从 ON 的中点垂直 ON 射入磁屏蔽区域，由几何关系可知 $r_{1} = L$

由洛伦兹力提供向心力

$qv_{0}B_{1} = m\frac{v_{0}^{2}}{r_{1}}$

联立可得 $B_{1} = \frac{mU}{qdB_{0}L}$

由于 $B_{2} < B_{1}$，根据洛伦兹力提供向心力

$qv_{0}B_{2} = m\frac{v_{0}^{2}}{r_{2}}$

解得 $r_{2} > L$

当 $B_{2} = 0$ 时粒子在磁屏蔽区向上做匀速直线运动，离开磁屏蔽区后根据左手定则，粒子向左偏转，如图所示。

\onepicture[h=5.50cm]{figs/14_an1.png}

根据洛伦兹力提供向心力

$qv_{0}B_{1} = m\frac{v_{0}^{2}}{r_{3}}$

可得 $r_{3} = r_{1} = L$

故粒子打在 $y$ 轴 $3L$ 处，综上所述 $y$ 轴上可能检测到该粒子的范围为 $L < y < 3L$。

（3）若在 Q 处检测到该粒子，如图。

\onepicture[h=5.50cm]{figs/14_an2.png}

由几何关系可知 $r_{2}^{2} = (2L)^{2} + (r_{2} - L)^{2}$

解得 $r_{2} = \frac{5}{2}L$

由洛伦兹力提供向心力

$qv_{0}B_{2} = m\frac{v_{0}^{2}}{r_{2}}$

联立解得 $B_{2} = \frac{2mU}{5qB_{0}dL}$

其中 $B_{1} = \frac{mU}{qdB_{0}L}$

根据磁屏蔽效率 $\eta = \frac{B_{1} - B_{2}}{B_{1}} \times 100\%$ 可得若在 Q 处检测到该粒子，则 $\eta = 60\%$。
}

\item
%% number: 15
%% paperName: 2025年云南高考第 15 题 15 分
%% typeId: 6
%% type: 计算题
%% chapter: 电磁感应
%% point: 电磁感应受力分析
%% method: 无
%% score: 15
%% degree: 350
%% duplicateId: 0
%% body:
如图所示，光滑水平面上有一个长为 $L$、宽为 $d$ 的长方体空绝缘箱，其四周紧固一电阻为 $R$ 的水平矩形导线框，箱子与导线框的总质量为 $M$。与箱子右侧壁平行的磁场边界平面如截面图中虚线 PQ 所示，边界右侧存在范围足够大的匀强磁场，其磁感应强度大小为 $B$、方向竖直向下。$t = 0$ 时刻，箱子在水平向右的恒力 $F$（大小未知）作用下由静止开始做匀加速直线运动，这时箱子左侧壁上距离箱底 $h$ 处、质量为 $m$ 的木块（视为质点）恰好能与箱子保持相对静止。箱子右侧壁进入磁场瞬间，木块与箱子分离；箱子完全进入磁场前某时刻，木块落到箱子底部，且箱子与木块均不反弹（木块下落过程中与箱子侧壁无碰撞）；木块落到箱子底部时即撤去 $F$。运动过程中，箱子右侧壁始终与磁场边界平行，忽略箱壁厚度、箱子形变、导线粗细及空气阻力。木块与箱子内壁间的动摩擦因数为 $\mu$，假设最大静摩擦力等于滑动摩擦力，重力加速度为 $g$。

\begin{enumerate}
\item
求 $F$ 的大小；
\item
求 $t = 0$ 时刻，箱子右侧壁距磁场边界的最小距离；
\item
若 $t = 0$ 时刻，箱子右侧壁距磁场边界的距离为 $s$（$s$ 大于（2）问中最小距离），求最终木块与箱子的速度大小。
\end{enumerate}

\twopicture[labelA=2025云南15a, labelB=2025云南15b, capA=绝缘箱立体图, capB=截面图, wA=6.20cm, wB=8.00cm, align=right]{figs/15a.png}{figs/15b.png}

%% answer:
\jdanswer{
\begin{enumerate}
\item
$\frac{(M + m)g}{\mu}$

\item
$\frac{(M + m)^{2}gR^{2}}{2\mu B^{4}d^{4}}$

\item
当 $\frac{g}{\mu}\left(\sqrt{\frac{2\mu s}{g}} + \sqrt{\frac{2h}{g}}\right) \geq \frac{B^{2}d^{2}L}{(M + m)R}$ 时，$v = \frac{g}{\mu}\left(\sqrt{\frac{2\mu s}{g}} + \sqrt{\frac{2h}{g}}\right) - \frac{B^{2}d^{2}L}{(M + m)R}$；当 $\frac{g}{\mu}\left(\sqrt{\frac{2\mu s}{g}} + \sqrt{\frac{2h}{g}}\right) < \frac{B^{2}d^{2}L}{(M + m)R}$ 时，$v = 0$。
\end{enumerate}
}
%% memo:
\memoanswer{
（1）对木块与箱子整体受力分析，由牛顿第二定律 $F = (M + m)a$

对木块受力分析，水平方向由牛顿第二定律 $F_{N} = ma$

竖直方向由平衡条件 $f = mg = \mu F_{N}$

联立可得 $F = \frac{(M + m)g}{\mu}$

（2）设箱子刚进入磁场中时速度为 $v$，产生的感应电动势为 $E = Bdv$

由闭合电路欧姆定律得，感应电流为 $I = \frac{E}{R}$

安培力大小为 $F_{\text{安}} = BId$

联立可得 $F_{\text{安}} = \frac{B^{2}d^{2}v}{R}$

若要使两物体分离，此时有 $F_{\text{安}} \geq F$

其中 $F = \frac{(M + m)g}{\mu}$

解得 $v \geq \frac{(M + m)gR}{\mu B^{2}d^{2}}$

由运动学公式 $v^{2} = 2as$

解得 $s \geq \frac{(M + m)^{2}gR^{2}}{2\mu B^{4}d^{4}}$

故 $t = 0$ 时刻，箱子右侧壁距磁场边界的最小距离为 $s_{\min} = \frac{(M + m)^{2}gR^{2}}{2\mu B^{4}d^{4}}$

（3）水平方向由运动学公式 $s = \frac{1}{2}at_{1}^{2}$

竖直方向有 $h = \frac{1}{2}gt_{2}^{2}$

其中 $F = \frac{(M + m)g}{\mu} = (M + m)a$

可得力 $F$ 作用的总时间为

$t = t_{1} + t_{2} = \sqrt{\frac{2\mu s}{g}} + \sqrt{\frac{2h}{g}}$

水平方向对系统由动量定理 $Ft - F_{\text{安}}\Delta t = (M + m)v - 0$

其中 $F_{\text{安}}\Delta t = \frac{B^{2}d^{2}L}{R}$

联立可得 $v = \frac{g}{\mu}\left(\sqrt{\frac{2\mu s}{g}} + \sqrt{\frac{2h}{g}}\right) - \frac{B^{2}d^{2}L}{(M + m)R}$

当 $\frac{g}{\mu}\left(\sqrt{\frac{2\mu s}{g}} + \sqrt{\frac{2h}{g}}\right) \geq \frac{B^{2}d^{2}L}{(M + m)R}$ 时，最终木块与箱子的速度大小为

$v = \frac{g}{\mu}\left(\sqrt{\frac{2\mu s}{g}} + \sqrt{\frac{2h}{g}}\right) - \frac{B^{2}d^{2}L}{(M + m)R}$

当 $\frac{g}{\mu}\left(\sqrt{\frac{2\mu s}{g}} + \sqrt{\frac{2h}{g}}\right) < \frac{B^{2}d^{2}L}{(M + m)R}$ 时，最终木块与箱子的速度大小为 $v = 0$。
}

\end{enumerate}

\end{document}
